\chapter{Protocol-Step Mapping Table}
\label{app:protocol-step-mapping}

The implementation uses function names that do not exactly match the step numbering of the
protocol paper. This appendix records the mapping used for all measurements.

\begin{table}[htbp]
  \centering
  \small
  \begin{tabularx}{\linewidth}{>{\raggedright\arraybackslash}p{0.27\linewidth}>{\raggedright\arraybackslash}p{0.18\linewidth}>{\raggedright\arraybackslash}X}
    \toprule
    Implementation action & Protocol step & Measurement meaning \\
    \midrule
    Contract creation / initialization & Step 1 & Deployment or clone creation of the
    optimistic exchange account. In the clone architecture, this is a marginal proxy
    creation once the implementation exists. \\
    \texttt{sendPayment} & Step 2 & Buyer payment. In the \(S=B\) optimized path this
    payment is fused with creation/initialization. \\
    \texttt{sendKey} & Step 3 & Vendor key release in the honest optimistic path. \\
    Buyer-side sponsor fee call & Step 4--5 & Buyer-side dispute sponsor
    registration and deposit. In the \(S_B=B\) case this is logically self-sponsored. \\
    Dispute trigger / vendor-side fee & Step 6 & Vendor-side dispute
    sponsor registration and deployment of the dispute account. In many tests this is a
    large transaction because it includes deployment. \\
    \texttt{respondChallenge} & Step 7b & Buyer publishes the current \(hpre_i\) value for
    the binary search. \\
    \texttt{giveOpinion} & Step 7c & Vendor accepts or rejects the buyer's response,
    updating the binary-search interval. \\
    Commitment submission & Step 8 & On-chain
    verification of the disputed gate, authenticated values, and accumulator proofs. \\
    Dispute finalization & Step 9 & Finalization or
    cancellation after the dispute outcome. \\
    \bottomrule
  \end{tabularx}
  \caption{Mapping between implementation actions and protocol steps.}
  \label{tab:app-step-mapping}
\end{table}

\section{Scope definitions}

\begin{table}[htbp]
  \centering
  \begin{tabularx}{\linewidth}{>{\raggedright\arraybackslash}p{0.28\linewidth}>{\raggedright\arraybackslash}X}
    \toprule
    Scope & Included operations \\
    \midrule
    Optimistic path & Step 1, Step 2, Step 3, and completion when no dispute occurs. \\
    Dispute trigger & Step 6, often including vendor-side sponsor deposit and deployment
    of the dispute account. \\
    Dispute after trigger & Step 7b, Step 7c, Step 8, and Step 9 after the dispute account
    exists. \\
    Complete dispute path & Optimistic setup, trigger, challenge rounds, Step 8, and
    finalization. \\
    Marginal clone cost & Per-exchange cost after factory and implementation contracts
    are already deployed. \\
    Fixed infrastructure cost & One-time deployment of reusable factory or implementation
    contracts. \\
    \bottomrule
  \end{tabularx}
  \caption{Gas-accounting scopes used throughout the report.}
  \label{tab:app-scope-definitions}
\end{table}

The most important caution is that \texttt{triggerDispute} should not be interpreted as a
small state transition. In the measured implementation, it often includes deployment of a
new dispute account, which explains why it dominates several complete dispute totals.
